\section{Limitations and Threats to Validity}
\label{sec:limitations}

The paper's strongest limitation is coverage. Although the registry contains 20
aliases, only seven were executed in the archived run, and those seven cover
only the Claude and Mistral families in practice. Any provider-level conclusion
outside that subset would therefore be overstated.

A second limitation is evaluation coupling. The archived \sr\ run used a single
judge model, \texttt{anthropic:claude-sonnet-4-6}. This provides consistency,
but it also leaves open the possibility of judge-specific bias. We do not have
human labels or a multi-judge ensemble for the current artifact set.

Third, the attack set is framework-native. The 20 ATT\&CK-inspired probes are a
useful operational taxonomy, but they are not the same as the public HarmBench
or \sr\ harmful-prompt datasets. Applying \sr\ to these probes is a deliberate
internal-evaluation choice, not a direct benchmark reproduction.

Fourth, the engine itself is not yet the empirical substrate of the paper. The
engine package is implemented, tested, and central to the framework's intended
design, but the archived results were generated with the stable comparison-runner
path. The paper therefore evaluates the framework through its current evidence
base rather than claiming a dedicated live engine sweep that did not happen.

Finally, the size figure is only exploratory. The reachable size-annotated
sample contains two Mistral models and no second family with multiple executed
size points. That is enough to motivate future work, but not enough to support a
general scaling narrative.
